\begin{table}[!h]
\centering
\caption{\label{tab:h6_robust}H6 robustness check: excluding ceiling cases (17 cases, 17.7\%).}
\centering
\fontsize{10}{12}\selectfont
\begin{threeparttable}
\begin{tabular}[t]{lcrrrr}
\toprule
Model & $n$ & $b$ & SE & OR & $p$\\
\midrule
\addlinespace[0.3em]
\multicolumn{6}{l}{\textbf{Continuous outcome (movement toward AI)}}\\
\hspace{1em}Full sample (LMM) & 96 & 1.23 & 0.33 & -- & $<$.001\\
\hspace{1em}Excluding ceiling cases (LMM) & 79 & 1.35 & 0.40 & -- & .001\\
\addlinespace[0.3em]
\multicolumn{6}{l}{\textbf{Binary outcome (moved toward AI: yes/no)}}\\
\hspace{1em}Full sample (GLMM) & 96 & 1.35 & 0.52 & 3.86 & .009\\
\hspace{1em}Excluding ceiling cases (GLMM) & 79 & 0.76 & 0.53 & 2.14 & .154\\
\bottomrule
\end{tabular}
\begin{tablenotes}
\item \textit{Note.} 
\item Ceiling cases are conversations where the participant could not move toward the AI because they were already at the extreme the AI argued for (e.g., rating = 7 when AI argued for high ratings). The H6 effect (opposing vs.\ reinforcing) remains significant in the continuous model but becomes non-significant in the binary model, suggesting some of the binary effect was driven by ceiling constraints.
\end{tablenotes}
\end{threeparttable}
\end{table}
